% Prostudováno: ucebnice/02_Funkce_a_rovnice/Linearni_funkce_a_rovnice/Linerání rovnice.tex,
% část Lineární rovnice o jedné neznámé, ekvivalentní úpravy a příklady (ř. 135–500).
% Realisticky.cz: externi_zdroje/realisticky/02_Funkce_a_rovnice/02_Linearni_rovnice_a_nerovnice/020204_Linearni_rovnice_II/020204_Lekce.pdf;
% využito pořadí úprav: nejprve závorky a jmenovatele, potom srovnání koeficientů. Úlohy vlastní.
\documentclass[11pt,a4paper]{article}
% Spolecny styl pro ucitelske pripravy MAT (Work/MAT)
% Verze 1.1, 2026-10-23
\NeedsTeXFormat{LaTeX2e}
\ProvidesPackage{priprava}[2026/10/23 v1.1 Sablona ucitelskych priprav MAT]

% Zakladni typografie a cestina
\RequirePackage[T1]{fontenc}
\RequirePackage[utf8]{inputenc}
\RequirePackage[czech]{babel}
\RequirePackage{lmodern}
\RequirePackage{microtype}
\RequirePackage{amsmath,amssymb}
\RequirePackage{xcolor}
\RequirePackage{enumitem}
\RequirePackage[most]{tcolorbox}
\RequirePackage{tabularx,booktabs,array}
\RequirePackage{fancyhdr}
\RequirePackage{needspace}
\RequirePackage{geometry}
\RequirePackage[hidelinks]{hyperref}

\geometry{a4paper,top=17mm,bottom=18mm,left=23mm,right=23mm}
\definecolor{notegray}{gray}{0.39}
\setlength{\parindent}{0pt}
\setlength{\parskip}{2.5pt plus 1pt minus 0.5pt}
\setlist[itemize]{topsep=2pt,itemsep=1pt,parsep=0pt,leftmargin=1.45em}
\setlist[enumerate]{topsep=2pt,itemsep=2pt,parsep=0pt,leftmargin=1.65em}
\setcounter{secnumdepth}{0}
\renewcommand{\arraystretch}{1.12}

% Zahlavi: vlevo rocnik | globalni tema, vpravo cislo: nazev hodiny
% Zapati zustava zcela prazdne.
\pagestyle{fancy}
\fancyhf{}
\renewcommand{\headrulewidth}{0.35pt}
\renewcommand{\footrulewidth}{0pt}
\setlength{\headheight}{14pt}
\newcommand{\nastavPripravu}[4]{%
  \fancyhead[L]{\small #1\enspace\textbar\enspace #2}%
  \fancyhead[R]{\small #3: #4}%
}

% Minimalisticke cernobile boxy. Kazda definice zvlast.
\tcbset{commonlesson/.style={enhanced,breakable=false,colback=white,colframe=black,
  colbacktitle=white,coltitle=black,boxrule=0.45pt,
  arc=0.8mm,outer arc=0.8mm,left=3mm,right=3mm,top=1.6mm,bottom=1.6mm,
  toptitle=1mm,bottomtitle=0.7mm,fonttitle=\bfseries,
  attach boxed title to top left={xshift=2.5mm,yshift=-1.15mm},
  boxed title style={colback=white,colframe=white,boxrule=0pt,left=1mm,right=1mm}}}
\newtcolorbox{zadaniBox}[1][]{commonlesson,title={Zadání},#1}
\newtcolorbox{prikladBox}[1][]{commonlesson,title={Příklad},#1}
\newtcolorbox{definiceBox}[1][]{commonlesson,title={Definice},#1}

% Spolecny box pro zadani a jeho reseni, bez mezery mezi castmi.
% Pouziti:
%   \begin{prikladReseniBox}
%     Zadani ulohy.
%     \tcblower
%     \begin{align*} ... \end{align*}
%   \end{prikladReseniBox}
% Spodni cast ma jemny sedy podklad a delici linku; lze v ni sazet
% vice radku matematiky. Obe casti maji jen jeden spolecny ramecek.
\tcbset{solutionlesson/.style={
  commonlesson,
  enhanced jigsaw,
  breakable,
  colbacklower=black!3,
  segmentation style={solid,draw=black!25,line width=0.35pt},
  fontlower=\small\color{notegray},
  before lower={\noindent\textbf{Řešení:}\enspace},
  leftlower=3mm,rightlower=3mm
}}
\newtcolorbox{prikladReseniBox}[1][]{solutionlesson,title={Příklad},#1}
\newtcolorbox{zadaniReseniBox}[1][]{solutionlesson,title={Zadání},#1}

% Pomucky pro autora pripravy: sedive poznamky; \reseni pro starsi pripravy
\newcommand{\poznamka}[1]{\par\begingroup\small\color{notegray}\itshape
  \textnormal{\textbf{Poznámka pro učitele.}}\enspace #1\par\endgroup}
\newcommand{\reseni}[1]{\par\begingroup\small\color{notegray}
  \textbf{Očekávané řešení:}\enspace #1\par\endgroup}
\newcommand{\blok}[1]{\needspace{5\baselineskip}\subsection*{#1}}
\newcommand{\vyrok}[1]{\uv{#1}}

\nastavPripravu{1. ročník}{Rovnice a nerovnice}{74}{Rovnice se závorkami a zlomky}
\begin{document}
\blok{Promyšlené pořadí úprav}
Z minulé hodiny víme, že stejná operace na obou stranách zachová řešení. U složitějšího zápisu se vyplatí nejdříve odstranit závorky a číselné jmenovatele, pak teprve sdružovat členy. Při násobení rovnice nejmenším společným násobkem jmenovatelů násobíme \emph{celé} obě strany včetně každého sčítance.
\begin{definiceBox}
\textbf{Společný násobek jmenovatelů} při řešení rovnice se zlomky je nenulové číslo dělitelné každým číselným jmenovatelem. Násobení obou stran takovým číslem je ekvivalentní a odstraní číselné zlomky.
\end{definiceBox}
\poznamka{V povinném jádru jsou pouze číselné jmenovatele, takže definiční obor je $\mathbb R$. Lomené výrazy s neznámou ve jmenovateli vyžadují podmínky a nepřidávej je bez výběru jako běžný krok.}
\begin{prikladReseniBox}
Řešte $3(x-2)-2(x+1)=4$. Proveďte zkoušku v původní rovnici.
\tcblower
Rozevřeme závorky: $3x-6-2x-2=4$, tedy $x-8=4$ a $x=12$. Zkouška: $3(12-2)-2(12+1)=30-26=4$. Množina řešení je $K=\{12\}$.
\end{prikladReseniBox}
\begin{prikladReseniBox}
Řešte $\dfrac{x-1}{3}+\dfrac{x+2}{2}=3$ a zdůvodněte násobení šesti.
\tcblower
Společný násobek jmenovatelů $3,2$ je $6\ne0$. Po násobení obou stran dostaneme $2(x-1)+3(x+2)=18$, tedy $2x-2+3x+6=18$, $5x=14$ a $x=\frac{14}{5}$. Zkouška: $(14/5-1)/3=3/5$ a $(14/5+2)/2=12/5$; součet je $15/5=3$. Proto $K=\{14/5\}$.
\end{prikladReseniBox}
\begin{prikladReseniBox}
Řešte $\dfrac{2x+1}{4}-\dfrac{x-2}{3}=1$. Vysvětlete znaménko před druhým zlomkem.
\tcblower
Násobíme $12$: $3(2x+1)-4(x-2)=12$. Mínus se vztahuje na celou závorku, proto $6x+3-4x+8=12$. Odtud $2x+11=12$ a $x=\frac12$. Zkouška: $(2\cdot\frac12+1)/4=\frac12$ a $(\frac12-2)/3=-\frac12$; rozdíl je $1$. Tedy $K=\{\frac12\}$.
\end{prikladReseniBox}
\blok{Samostatné procvičení}
\begin{zadaniReseniBox}
\begin{enumerate}[label=\alph*)]
\item Řešte $4(2x-1)-3(x+2)=5$.
\item Řešte $\dfrac{x+1}{2}=\dfrac{x-2}{3}+2$.
\item Rozhodněte o počtu řešení $2(x+4)=2x+8$ a $2(x+4)=2x+7$.
\end{enumerate}
\tcblower
\begin{enumerate}[label=\alph*)]
\item $8x-4-3x-6=5$, $5x=15$, $x=3$. Kontrola: $4(5)-3(5)=5$.
\item Násobíme $6$: $3(x+1)=2(x-2)+12$, tedy $3x+3=2x+8$ a $x=5$. Kontrola: $6/2=3$ a $3/3+2=3$.
\item První je po roznásobení identita $2x+8=2x+8$, takže $K=\mathbb R$. Druhá vede k $8=7$, tedy $K=\varnothing$.
\end{enumerate}
\end{zadaniReseniBox}
\blok{Diagnostika chyb}
\begin{zadaniReseniBox}
Žák napsal z $\frac{x+1}{2}+\frac{x-1}{3}=4$ po násobení šesti rovnici $3(x+1)+(x-1)=24$. Opravte ji a rovnici dopočítejte.
\tcblower
Druhý zlomek se po násobení šesti násobí $6/3=2$, ne jedničkou. Správně $3(x+1)+2(x-1)=24$, tedy $3x+3+2x-2=24$, $5x=23$, $x=23/5$. Kontrola: $(28/5)/2=14/5$, $(18/5)/3=6/5$, součet $4$.
\end{zadaniReseniBox}
\blok{Volitelné rozšíření a rezerva}
\begin{zadaniReseniBox}
Řešte $\dfrac{2(x-1)-3}{5}=\dfrac{x+4}{2}-2$ a prověřte, zda nenastala ztráta řešení.
\tcblower
Násobíme $10$: $2(2x-5)=5(x+4)-20$. Odtud $4x-10=5x$, tedy $x=-10$. Zkouška: vlevo $(2(-11)-3)/5=-25/5=-5$, vpravo $(-6)/2-2=-5$. Číselné jmenovatele $5,2$ jsou nenulové; ztráta řešení nenastala.
\end{zadaniReseniBox}
\blok{Domácí práce}
\begin{zadaniReseniBox}
\begin{enumerate}[label=\alph*)]
\item Řešte $5(x-2)-2(2x+1)=3$ a proveďte zkoušku.
\item Řešte $\dfrac{x-3}{4}+\dfrac{x+1}{2}=1$.
\item Vysvětlete rozdíl mezi násobením rovnice číslem $6$ a násobením jen jednoho jejího členu.
\end{enumerate}
\tcblower
\begin{enumerate}[label=\alph*)]
\item $5x-10-4x-2=3$, tedy $x=15$; kontrola $5(13)-2(31)=65-62=3$.
\item Násobíme $4$: $x-3+2(x+1)=4$, tedy $3x-1=4$ a $x=5/3$. Kontrola: $(5/3-3)/4=-1/3$, $(5/3+1)/2=4/3$, součet $1$.
\item Ekvivalentní je násobit \emph{obě celé strany} stejným nenulovým číslem. Násobení jediného členu mění hodnotu strany, a zpravidla tedy i množinu kořenů.
\end{enumerate}
\end{zadaniReseniBox}
\poznamka{Jádro: strategie 5 min, společné příklady 17 min, samostatná práce včetně kontroly 20 min, organizace 3 min. Diagnostiku využij při opravě; rezerva nad rámec. Domácí práce 15–20 min. Před další hodinou upevni zápis intervalů.}
\end{document}
